\documentclass[aspectratio=169]{beamer}
\usetheme{Madrid}
\usecolortheme{default}
\usepackage{booktabs}
\usepackage{graphicx}
\usepackage{amsmath}
\usepackage{microtype}

% Suppress navigation bar
\setbeamertemplate{navigation symbols}{}
\setbeamertemplate{footline}[frame number]

% Figure path
\graphicspath{{../figures/}}

\title{Seasonal Employment: Replication and Extension}
\subtitle{Coglianese \& Price (2025) + QWI 6-Digit NAICS Index}
\date{\today}

\begin{document}

% ── Title ─────────────────────────────────────────────────────────────────────
\begin{frame}
  \titlepage
\end{frame}

% ── Outline ───────────────────────────────────────────────────────────────────
\begin{frame}{Outline}
  \tableofcontents
\end{frame}

% ══════════════════════════════════════════════════════════════════════════════
\section{Replication of Coglianese \& Price (2025)}
% ══════════════════════════════════════════════════════════════════════════════

\begin{frame}{Coglianese \& Price (2025): Overview}
  \begin{itemize}
    \item ``Income in the Off-Season: Household Adaptation to Yearly Work
          Interruptions,'' \textit{ILRR} 2025
    \item \textbf{Question:} How do households smooth income when employment
          separations recur annually (i.e.\ are \emph{seasonal})?
    \item \textbf{Key innovation:} Identify seasonal separations from
          administrative data without requiring industry or occupation labels
    \bigskip
    \item Core estimand: \textbf{excess recurrence} at 12-month horizon
    \[
      \widehat{\text{ExR}}
      = \Pr(\text{sep at }k{=}12)
        - \frac{\Pr(\text{sep at }k{=}11) + \Pr(\text{sep at }k{=}13)}{2}
    \]
    \item Counterfactual = average of adjacent lags $\to$ isolates the
          \emph{annual spike} marking genuine seasonal separations
  \end{itemize}
\end{frame}

\begin{frame}{Estimation Strategy}
  \textbf{Event-study OLS on separators panel (CPS and SIPP):}
  \[
    \mathbf{1}[\text{sep at }k] = \sum_{j} \beta_j \mathbf{1}[k=j] + \varepsilon
  \]
  \begin{itemize}
    \item Estimate by OLS; excess recurrence = $\beta_{12} - (\beta_{11}+\beta_{13})/2$
    \item Standard error via delta method
    \bigskip
    \item \textbf{Two specifications:}
    \begin{enumerate}
      \item \textit{Raw} (Figure 2 in paper): only horizon dummies, no constant
      \item \textit{Preferred} (Appendix F.1): adds controls for interview-seam
            artifacts (SIPP) and weeks elapsed since separation (CPS)
    \end{enumerate}
    \bigskip
    \item Python replication uses \texttt{statsmodels} WLS with clustered SEs;
          data from the LEHD-style pre-cleaned \texttt{.dta.gz} files in the
          replication package
  \end{itemize}
\end{frame}

\begin{frame}{Replication Results}
  \begin{table}
    \centering
    \caption{Excess recurrence estimates: paper vs.\ Python replication}
    \begin{tabular}{llcc}
      \toprule
      Specification & Where & Paper & Replication \\
      \midrule
      CPS raw (Figure 2) & Figure 2 & 1.4 p.p.\ (SE 0.14) & 1.44\ \checkmark \\
      CPS preferred & App.\ F.1 & 1.5 p.p.\ (SE 0.14) & 1.42\ \checkmark \\
      SIPP raw (Figure 2) & Figure 2 & 2.0 p.p.\ (SE 0.11) & 2.04\ \checkmark \\
      SIPP preferred & App.\ F.1 & 1.6 p.p.\ (SE 0.11) & 1.62\ \checkmark \\
      \bottomrule
    \end{tabular}
  \end{table}
  \vspace{0.5em}
  \begin{itemize}
    \item All four estimates within rounding of published values
    \item Industry ranking (both datasets): Agriculture $>$ Entertainment
          $>$ Education $>$ Construction $\gg$ Healthcare, FIRE
  \end{itemize}
\end{frame}

% ══════════════════════════════════════════════════════════════════════════════
\section{QWI 6-Digit NAICS Seasonal Index}
% ══════════════════════════════════════════════════════════════════════════════

\begin{frame}{Extending to 6-Digit NAICS via QWI}
  \begin{itemize}
    \item CP's classifier operates at the worker level; industry rankings are
          reported at 1-digit NAICS only
    \item \textbf{Goal:} produce a finer, industry-level seasonal index suitable
          for firm classification
    \bigskip
    \item \textbf{Data:} Census Quarterly Workforce Indicators (QWI)
    \begin{itemize}
      \item Quarterly separations and employment by state $\times$ 6-digit NAICS
      \item 51 states + DC, years 2000--2023 (1,011 industries)
    \end{itemize}
    \bigskip
    \item \textbf{Quarterly analog to excess recurrence:} for each quarter
          $q \in \{1,2,3,4\}$ and industry $j$,
    \[
      \text{Excess}(q)_{j}
      = \overline{
          \text{SepRate}_{j,q,t}
          - \frac{\text{SepRate}_{j,q-1,t} + \text{SepRate}_{j,q+1,t}}{2}
        }
    \]
    (quarters treated cyclically; average over years $t$)
    \item \textbf{Peak excess:} $\text{ExcessPeak}_{j} = \max_q\, \text{Excess}(q)_j$
    \item \textbf{Seasonal amplitude:} $\max_q - \min_q$ normalised to $[0,1]$;
          does not assume any particular quarter is ``the'' seasonal quarter
  \end{itemize}
\end{frame}

\begin{frame}{National Seasonal Index: Top Industries}
  \begin{columns}
    \begin{column}{0.55\textwidth}
      \begin{table}
        \footnotesize
        \caption{Top 12 most seasonal 6-digit NAICS industries (national)}
        \begin{tabular}{lcc}
          \toprule
          NAICS \& description & ExcessPeak & $n$ \\
          \midrule
          111336 Fruit/tree nut combo farm.  & 0.333 & 18 \\
          722330 Mobile food services        & 0.316 & 18 \\
          111140 Wheat farming                & 0.287 & 20 \\
          111160 Rice farming                 & 0.561 & 22 \\
          111199 Other grain farming          & 0.397 & 23 \\
          111219 Other veg.\ (ex. potato) farm.\ & 0.459 & 15 \\
          111998 Other misc.\ crop farming     & 0.328 & 20 \\
          111120 Oilseed (ex.\ soybean) farm.\ & 0.347 & 12 \\
          115113 Crop harvesting (machine)     & 0.398 & 10 \\
          512132 Drive-in motion picture theaters & 0.362 & 21 \\
          532284 Recreational goods rental     & 0.485 & 22 \\
          111191 Oilseed/grain combo farm.\ & 0.321 & 23 \\
          \bottomrule
        \end{tabular}
      \end{table}
    \end{column}
    \begin{column}{0.42\textwidth}
      \begin{itemize}
        \item $n$ = years contributing to ExcessPeak estimate; a $\geq 5$-year
              minimum is enforced so thin, disclosure-suppressed cells can't
              masquerade as extreme seasonality
        \item Non-agriculture entries are intuitive once named: mobile food
              services (food trucks), drive-in theaters, recreational
              goods rental
        \item Seasonal index normalized to [0,1] via 99th-pct winsorization
        \item Correlation with CP 1-digit rankings: $\rho = 0.835$
      \end{itemize}
    \end{column}
  \end{columns}
\end{frame}

\begin{frame}{National Seasonal Index: Sector Summary}
  \begin{table}
    \centering
    \footnotesize
    \caption{Mean seasonal index by 2-digit sector (national)}
    \begin{tabular}{lcc}
      \toprule
      Sector & Mean seasonal index & ExcessPeak (p.p.) \\
      \midrule
      Agriculture, forestry, fishing      & 0.471 & 15.0 \\
      Arts, entertainment \& recreation   & 0.445 & 12.1 \\
      Accommodation \& food services      & 0.323 & 8.3  \\
      Educational services                & 0.273 & 8.4  \\
      Construction                        & 0.194 & 5.9  \\
      Real estate                         & 0.167 & 5.6  \\
      Warehousing \& couriers             & 0.135 & 4.6  \\
      Other services                      & 0.134 & 4.2  \\
      \midrule
      Health care \& social assistance    & 0.067 & 2.2  \\
      Wholesale trade                     & 0.063 & 1.8  \\
      Manufacturing (metals/machinery)    & 0.036 & 1.1  \\
      \bottomrule
    \end{tabular}
  \end{table}
  \vspace{0.3em}
  \small Education peaks in Q2 (end of school year), not Q4 --- ExcessPeak is
  positive across all sectors. CP's 1-digit rankings confirmed at finer granularity.
\end{frame}

\begin{frame}{Robustness Check: Does COVID Distort the Index?}
  \begin{itemize}
    \item The national index pools 2000--2023. Mass pandemic-era layoffs in
          industries already ranked highly seasonal (accommodation/food,
          arts \& entertainment) could inflate ``excess separation'' for
          reasons that have nothing to do with normal annual seasonality
    \item \textbf{Test:} recompute the index restricted to 2000--2019 only,
          compare to the full sample
    \bigskip
    \item \textbf{Result: rankings are robust.}
    \begin{itemize}
      \item Pearson correlation (full vs.\ pre-COVID seasonal index): 0.997
      \item Spearman rank correlation: 0.986
      \item Only 6.3\% of industries have a different peak quarter
      \item Watch sectors (agriculture, construction, education, arts/entertainment,
            accommodation/food) all shift by $\leq 0.017$
    \end{itemize}
    \bigskip
    \item \textbf{One genuine mover:} a transportation industry (transit,
          NAICS 485111) whose index is roughly 5$\times$ higher with COVID years
          included --- plausibly a real ridership-collapse effect, not a data artifact
    \item[$\to$] Running this check is what surfaced a stale-data-threshold bug:
          three industries (incl.\ Tax Preparation Services, Cotton Ginning) had
          been wrongly ranked \#1 nationally off just 2 years of COVID-distorted
          data, below the 5-year minimum --- now fixed at the source
  \end{itemize}
\end{frame}

% ══════════════════════════════════════════════════════════════════════════════
\section{Has Seasonality Changed Over Time?}
% ══════════════════════════════════════════════════════════════════════════════

\begin{frame}{Seasonality Over Time: 2000--2010 vs.\ 2011--2023}
  \begin{itemize}
    \item The national index pools all 24 years. Has seasonality actually been
          stable, or has it drifted over two decades?
    \item \textbf{Test:} split the sample in two and recompute independently.
          The 5-year $n\_excess\_obs$ minimum is what makes a clean two-way
          split possible --- both halves (11 and 13 years) keep near-full
          coverage, so a three- or four-way split isn't advisable
    \bigskip
    \item \textbf{Result: correlated, with real drift.}
    \begin{itemize}
      \item Pearson $r = 0.869$, Spearman $r = 0.795$ (both notably lower
            than the COVID check's 0.997/0.986 --- an 11-year gap between
            period midpoints is a much bigger ask than excluding 2 years)
      \item Economy-wide mean seasonal index is essentially flat:
            $0.130 \to 0.126$
    \end{itemize}
    \bigskip
    \item \textbf{But the flat average hides a real reshuffling} beneath it
          (next slide)
    \item[$\to$] 30.5\% of industries have a different peak quarter, but this
          concentrates in weakly-seasonal industries (44\% of the bottom
          quartile vs.\ 15\% of the top) --- ordinary noise, not genuine
          timing instability among strongly seasonal industries
  \end{itemize}
\end{frame}

\begin{frame}{Which Sectors Got More or Less Seasonal?}
  \begin{columns}
    \begin{column}{0.55\textwidth}
      \begin{center}
        \includegraphics[width=\textwidth,height=0.8\textheight,keepaspectratio]%
          {seasonality_trend_by_sector.pdf}
      \end{center}
    \end{column}
    \begin{column}{0.4\textwidth}
      \scriptsize
      \textbf{Agriculture: more seasonal}
      \begin{itemize}
        \item $+0.044$, the largest sector increase (0.418 $\to$ 0.462)
        \item Plausibly consistent with growth in H-2A seasonal
              guest-worker visa usage over the 2010s (workers who leave
              every season, not year-round staff) --- a hypothesis, not
              independently verified here
      \end{itemize}
      \smallskip
      \textbf{Construction: less seasonal}
      \begin{itemize}
        \item $-0.046$, the largest sector decrease (0.216 $\to$ 0.170) ---
              the single biggest shift either direction
      \end{itemize}
      \smallskip
      \textbf{Most others: mild decline}
      \begin{itemize}
        \item Accommodation/food, transportation, professional services,
              arts/entertainment all $-0.01$ to $-0.03$
      \end{itemize}
    \end{column}
  \end{columns}
\end{frame}

% ══════════════════════════════════════════════════════════════════════════════
\section{Geographic Variation}
% ══════════════════════════════════════════════════════════════════════════════

\begin{frame}{Geographic Variation: Overview}
  \begin{itemize}
    \item The national index masks substantial state-level heterogeneity
    \item We compute ExcessPeak for each state $\times$ 6-digit NAICS cell
          (40,222 cells with $\geq 5$ year-observations)
    \bigskip
    \item \textbf{Main finding:} a 2.4$\times$ gap between most and least seasonal
          states in mean peak excess (Alaska: 8.0 p.p.\ vs.\ Texas: 3.3 p.p.)
    \item Northern/Great Plains states consistently more seasonal across
          multiple sectors; Sun Belt states consistently less so
    \bigskip
    \item \textbf{Most geographically variable sectors:}
    \begin{enumerate}
      \item Agriculture (cross-state SD = 15.4 p.p.)
      \item Arts \& entertainment (SD = 14.4 p.p.)
      \item Construction (SD = 9.9 p.p.)
    \end{enumerate}
    \item \textbf{Least geographically variable:} healthcare, metals manufacturing,
          utilities --- consistent low seasonality everywhere
  \end{itemize}
\end{frame}

\begin{frame}{State-Level Mean Seasonality}
  \begin{center}
    \includegraphics[width=0.95\textwidth,height=0.78\textheight,keepaspectratio]%
      {geo_state_mean_tile_map.png}
  \end{center}
\end{frame}

\begin{frame}{Construction and Agriculture by State}
  \begin{center}
    \includegraphics[width=0.95\textwidth,height=0.68\textheight,keepaspectratio]%
      {geo_construction_agriculture_by_state.pdf}
  \end{center}
  \vspace{-0.3em}
  {\small
  \textbf{Construction:} Minnesota (19.7 p.p.) vs.\ Florida (2.6 p.p.) --- an 8$\times$
  gap driven by winter construction shutdowns in northern states.
  \textbf{Agriculture:} Montana (20.2 p.p.) and Idaho (19.3 p.p.) lead.}
\end{frame}

\begin{frame}{State-Specific Top and Bottom 1\% Industries}
  \begin{columns}
    \begin{column}{0.50\textwidth}
      \begin{table}
        \footnotesize
        \caption{Within-state NAICS6 tails}
        \begin{tabular}{lcc}
          \toprule
          Statistic & Top 1\% & Bottom 1\% \\
          \midrule
          State-industry slots & 413 & 413 \\
          Unique NAICS6 codes & 133 & 252 \\
          Codes appearing once & 57 & 158 \\
          Codes in 10+ states & 6 & 0 \\
          Pairwise Jaccard & 0.054 & 0.015 \\
          Also national same tail & 10.9\% & 2.7\% \\
          \bottomrule
        \end{tabular}
      \end{table}
      \vspace{0.4em}
      {\small
      Interpretation: the state-specific tails are related to the national index,
      but most state-tail industries are not simply the national top or bottom
      1\% repeated everywhere.}
    \end{column}
    \begin{column}{0.46\textwidth}
      \footnotesize
      \textbf{Recurring top-1\% industries}
      \begin{itemize}
        \item Nursery/tree production: 15 states
        \item Marinas: 14 states
        \item Specialty trade contractors: 12 states
        \item Golf courses/country clubs: 11 states
        \item Miscellaneous crop farming: 11 states
      \end{itemize}
      \medskip
      \textbf{Bottom-1\% pattern}
      \begin{itemize}
        \item More diffuse than the top tail
        \item Concentrated in metals manufacturing, wholesale,
              chemicals, finance, and transportation
      \end{itemize}
    \end{column}
  \end{columns}
\end{frame}

\begin{frame}{Geographic Variation: Sector Cross-State SD}
  \begin{center}
    \includegraphics[width=0.85\textwidth,height=0.78\textheight,keepaspectratio]%
      {geo_sector_variation.pdf}
  \end{center}
\end{frame}

% ══════════════════════════════════════════════════════════════════════════════
\section{County-Level Exposure (California Pilot)}
% ══════════════════════════════════════════════════════════════════════════════

\begin{frame}{County-Level Seasonal Exposure: Methodology}
  \begin{itemize}
    \item \textbf{Goal:} how seasonal is the average job in a given county,
          given the industries actually located there?
    \[
      \text{Exposure}_c = \frac{\sum_j \text{EMP}_{c,j} \times \text{seasonal\_index}_j}
                               {\sum_j \text{EMP}_{c,j}}
    \]
    \item \textbf{Data:} County Business Patterns (CBP) establishment/employment
          counts by county $\times$ NAICS, merged onto the state $\times$
          NAICS6 seasonal index
    \bigskip
    \item \textbf{Problem:} CBP suppresses most county $\times$ NAICS6 cells
          (disclosure avoidance) --- a county's genuinely seasonal niche
          industries are often exactly the small, easily-identified cells
          that get suppressed
    \item \textbf{Fix: three-tier fallback} for each industry's seasonal
          score: (1) state $\times$ NAICS6, (2) national NAICS6, (3) national
          NAICS4 (also recovers employment CBP suppresses at 6-digit but
          still reports at 4-digit)
    \item[$\to$] California pilot: \textbf{100\% employment coverage} in all
          58 counties --- no county's score rests on a partial industry match
  \end{itemize}
\end{frame}

\begin{frame}{California County Ranking}
  \begin{center}
    \includegraphics[width=0.85\textwidth,height=0.82\textheight,keepaspectratio]%
      {county_seasonal_exposure_06.pdf}
  \end{center}
\end{frame}

\begin{frame}{Case Study: Yolo County}
  \begin{itemize}
    \item \textbf{Yolo County ranks 40th of 58} California counties (32nd
          percentile) --- \emph{less} exposed to seasonal work than roughly
          two-thirds of the state, despite being an agricultural county
    \bigskip
    \item \textbf{Why:} employment-weighting means large, stable
          service/logistics employers dominate the score:
    \begin{itemize}
      \item Limited-service restaurants (3,222 jobs, index 0.10), general
            warehousing (3,087 jobs, index 0.06), couriers (1,731 jobs,
            index 0.14) --- all low-seasonality but large employers
      \item Crop harvesting scores the \emph{maximum} 1.00 --- but employs
            only 246 people locally
    \end{itemize}
    \bigskip
    \item \textbf{Context:} Central Valley farm counties top the ranking
          (Madera 0.179, Colusa 0.170 --- roughly 2$\times$ Yolo's 0.088);
          the rest of the Sacramento metro area (Sacramento, Solano, Placer)
          clusters near Yolo at the low-seasonality end
    \item[$\to$] A county can have extreme seasonal niches and still score
          low overall, if its broader economy has diversified around them
  \end{itemize}
\end{frame}

% ══════════════════════════════════════════════════════════════════════════════
\section{Flow vs.\ Stock: Two Views of Seasonality}
% ══════════════════════════════════════════════════════════════════════════════

\begin{frame}{Flow vs.\ Stock: Motivation}
  \begin{itemize}
    \item The index has always been built two ways in parallel, but the two
          were never formally compared:
    \begin{itemize}
      \item \textbf{Flow} (\texttt{seasonal\_index}): from separation
            \emph{rates} --- does a quarter see an excess of separations?
      \item \textbf{Stock} (\texttt{seasonal\_index\_emp}): from employment
            \emph{levels} normalized to their own annual mean --- does
            headcount itself swell or shrink that quarter?
    \end{itemize}
    \bigskip
    \item \textbf{Result: correlated, but far from the same thing.}
    \begin{itemize}
      \item National ($n=1{,}004$): Pearson $r = 0.653$, Spearman $r = 0.627$
      \item State $\times$ NAICS6 ($n=39{,}225$): Pearson $r = 0.521$,
            Spearman $r = 0.585$
      \item \textbf{Peak-quarter concordance is low --- only $\sim$15\%
            nationally, $\sim$14\% by state}: even when both measures agree
            an industry is seasonal, they usually disagree on \emph{which
            quarter} drives it
    \end{itemize}
    \item[$\to$] Timing, not just amplitude, differs between ``separations
          spike'' and ``headcount swings''
  \end{itemize}
\end{frame}

\begin{frame}{Flow vs.\ Stock: Sector Heterogeneity}
  \begin{columns}
    \begin{column}{0.52\textwidth}
      \begin{center}
        \includegraphics[width=\textwidth,height=0.72\textheight,keepaspectratio]%
          {flow_vs_stock_national_scatter.pdf}
      \end{center}
    \end{column}
    \begin{column}{0.44\textwidth}
      \scriptsize
      \textbf{Construction: near-perfect agreement}
      \begin{itemize}
        \item Flow/stock correlation $r \approx 0.96$--$0.97$
        \item A direct, physical relationship: winter layoffs $=$ winter headcount drop
      \end{itemize}
      \smallskip
      \textbf{Agriculture: the opposite}
      \begin{itemize}
        \item High seasonal index on \emph{both} measures, but $r \approx 0.21$
              between them
        \item Consistent with rapid within-quarter worker replacement: departing
              workers are immediately replaced, spiking separations without
              moving headcount much
      \end{itemize}
      \smallskip
      \textbf{Validates the bug fix}
      \begin{itemize}
        \item Tax Prep Services and Cotton Ginning --- wrongly inflated on the
              flow side (previous slide) --- now correctly show \emph{low}
              flow, \emph{high} stock: a real hiring-surge story, not an artifact
      \end{itemize}
    \end{column}
  \end{columns}
\end{frame}

% ══════════════════════════════════════════════════════════════════════════════
\section{Income/Earnings Seasonality}
% ══════════════════════════════════════════════════════════════════════════════

\begin{frame}{A Third Lens: Does Income Dip in the Off-Season?}
  \begin{itemize}
    \item Flow and stock both measure separations/headcount. A third QWI
          measure, \texttt{EarnBeg} (average monthly earnings), tests the
          paper's title question directly: \textbf{does income dip in an
          industry's own off-season?}
    \item \textit{Data note:} the originally-requested \texttt{Payroll}
          field returns \texttt{null} from the Census API for every cell
          tested (even huge industries where employment populates fine);
          \texttt{EarnBeg} is the working substitute
    \bigskip
    \item \textbf{A confound has to be fixed first:} 83\% of \emph{all}
          1,011 industries --- not just seasonal ones --- have their raw
          earnings index peak in Q4. That's a broad, economy-wide
          year-end bonus/holiday-pay effect, not industry-specific timing
    \item Left uncorrected, this swamps the signal: whether an industry's
          \emph{own} separation-peak quarter happens to land in Q4 (everyone's
          bonus quarter) almost entirely determines whether it looks like
          income "dips" there
    \item[$\to$] \textbf{Fix:} compute the common cross-industry calendar
          effect (Q1 $-0.031$, Q2 $-0.005$, Q3 $-0.048$, Q4 $+0.085$) and net
          it out before testing --- the \emph{idiosyncratic} index isolates
          industry-specific timing
  \end{itemize}
\end{frame}

\begin{frame}{Does Income Dip in the Off-Season? The Answer}
  \begin{columns}
    \begin{column}{0.52\textwidth}
      \begin{center}
        \includegraphics[width=\textwidth,height=0.72\textheight,keepaspectratio]%
          {earnings_at_separation_peak_hist.pdf}
      \end{center}
    \end{column}
    \begin{column}{0.44\textwidth}
      \scriptsize
      \textbf{Naive test: 51.9\% "dip"} --- but almost entirely mechanical:
      \begin{itemize}
        \item Only 8.3\% dip when the separation peak falls in Q4 (everyone's
              earnings peak)
        \item 93.6\% "dip" when it falls in Q3 (rarely anyone's earnings peak)
      \end{itemize}
      \smallskip
      \textbf{De-confounded (credible) test:}
      \begin{itemize}
        \item \textbf{43.9\% dip}, mean excess $+0.0065$ --- essentially
              flat, not clearly negative
        \item No strong evidence of industry-average income dipping in the
              off-season
      \end{itemize}
      \smallskip
      \textbf{Construction: the sharp counter-example}
      \begin{itemize}
        \item Dip rate only 3.2\%, mean excess $+0.044$ --- earnings
              \emph{higher}, plausibly overtime pay before winter layoffs
      \end{itemize}
      \smallskip
      \textit{Doesn't contradict CP: their effect is the separating worker's
      own income loss (micro); this tests the industry-wide average (macro)
      --- a different, complementary question.}
    \end{column}
  \end{columns}
\end{frame}

% ══════════════════════════════════════════════════════════════════════════════
\section{Discussion}
% ══════════════════════════════════════════════════════════════════════════════

\begin{frame}{Implications for Firm Classification}
  \begin{itemize}
    \item The QWI index enables firm-level seasonal classification via NAICS lookup
    \item The top/bottom 1\% analysis shows the exact NAICS6 tails differ sharply
          across states: top tail has 133 unique codes; bottom tail has 252
    \item \textbf{But:} for geographically sensitive sectors (construction,
          agriculture), a single national index may misclassify:
    \begin{itemize}
      \item A Minnesota construction firm is highly seasonal
      \item A Florida construction firm is nearly non-seasonal
    \end{itemize}
    \bigskip
    \item \textbf{Preferred approach for firm classification:}
    \begin{enumerate}
      \item Use the \emph{state $\times$ NAICS} ExcessPeak as the lookup value
            when state of operation is known
      \item Fall back to national index for industries with thin state-level
            coverage (low $n\_excess\_obs$)
      \item Direct estimation from UI wage records when sufficient separations
            are available ($n \geq 20$ per firm)
    \end{enumerate}
    \smallskip
    \item \textbf{Which index?} Flow suits classification by \emph{turnover
          risk}; stock suits \emph{headcount/staffing swings} --- they agree
          only $\sim$65\% of the time nationally
    \item \textbf{Caveat:} 6-digit cells are often suppressed in small states;
          state-level estimates for niche agriculture sub-industries may rest on
          very few year-observations
  \end{itemize}
\end{frame}

\begin{frame}{Summary}
  \scriptsize
  \begin{enumerate}
    \item \textbf{Replication:} Python replication of CP (2025) excess recurrence
          matches all four published estimates within rounding
    \smallskip
    \item \textbf{National index:} QWI-based ExcessPeak at 6-digit NAICS correlates
          strongly with CP's 1-digit rankings ($\rho = 0.835$); robust to
          excluding 2020--2021 (Pearson $0.997$). Economy-wide seasonality is
          roughly flat 2000--10 to 2011--23, but Agriculture rose ($+0.044$)
          and Construction fell ($-0.046$) beneath that average
    \smallskip
    \item \textbf{Geographic variation:} large state heterogeneity (2.4$\times$
          gap, Alaska vs.\ Texas); a California county-level pilot applies the
          same logic down to counties --- Yolo County ranks 40th of 58 CA
          counties, less seasonal than most despite its agricultural economy
    \smallskip
    \item \textbf{Flow vs.\ stock:} separations-based and employment-based
          seasonality are correlated but distinct ($r \approx 0.65$ national);
          sharp sector heterogeneity (Construction: near-identical; Agriculture: not)
    \smallskip
    \item \textbf{Income/earnings:} a third lens (\texttt{EarnBeg}) finds no
          strong evidence of industry-average income dipping in the off-season
          once a common Q4-bonus effect is removed; Construction again the
          outlier (earnings \emph{higher}, not lower)
    \smallskip
    \item \textbf{Next steps:} national CBP rollout for a full county-level
          index; apply the state $\times$ NAICS index to UI wage records for
          firm-level classification
  \end{enumerate}
\end{frame}

\end{document}
